\documentclass[12pt]{article}
\usepackage[utf8]{inputenc}
\usepackage[margin=1in]{geometry}
\usepackage{graphicx}
\usepackage{booktabs}
\usepackage{amsmath}
\usepackage{hyperref}
\usepackage{natbib}
\usepackage{xcolor}
\usepackage{multirow}

\title{Deciphering Target Specificity and Inhibition Mechanisms of Pyrin-Only Proteins in Inflammasome Regulation}
\author{~}
\date{}

\begin{document}
\maketitle

% ============================================================
\begin{abstract}
Inflammasomes are multi-protein signalling platforms that activate inflammatory caspases in response to pathogen- and danger-associated molecular patterns. Dysregulated inflammasome activity underlies autoinflammatory diseases, yet the endogenous mechanisms that restrain these complexes remain incompletely understood. Pyrin-only proteins (POP1 and POP2) are small, single-domain regulators hypothesised to inhibit inflammasome assembly through homotypic pyrin domain (PYD) interactions. Here we combine Rosetta-based computational modelling, in vitro filament reconstitution assays, and cell-based inflammasome activation experiments to comprehensively map the target specificity and inhibition mechanisms of POP1 and POP2. Contrary to a previous report, we find that POP1 is a poor inhibitor of the central adaptor ASC; instead, POP1 potently blocks upstream receptor PYD filament assembly for AIM2, IFI16, NLRP3, and NLRP6. POP2 operates through a distinct dual mechanism: it directly suppresses ASC nucleation and additionally caps elongating AIM2 and NLRP6 filaments. These findings establish that POP1 and POP2 employ non-redundant strategies to regulate inflammasome signalling at different nodes of the pathway and provide design principles for engineering selective inflammasome inhibitors.
\end{abstract}

% ============================================================
\section{Introduction}

Inflammasomes are cytoplasmic multi-protein complexes that serve as sentinel platforms of the innate immune system \citep{martinon2002inflammasome}. Upon detecting pathogen-associated molecular patterns (PAMPs) or danger-associated molecular patterns (DAMPs), sensor proteins such as NLRP3, AIM2, and NLRP6 oligomerise through their pyrin domains (PYDs), nucleating helical filaments of the adaptor protein ASC, which in turn recruits and activates inflammatory caspases \citep{lu2014unified, broz2016inflammasomes}.

While inflammasome activation is essential for host defence, excessive or chronic activation drives a spectrum of autoinflammatory diseases, including cryopyrin-associated periodic syndromes, familial Mediterranean fever, gout, and contributes to cancer progression \citep{masters2009horror, karki2017diverging}. Consequently, understanding the endogenous mechanisms that restrain inflammasome activity is a pressing goal for both basic biology and therapeutic development.

Pyrin-only proteins (POPs) are small ($\sim$10\,kDa), single-PYD proteins expressed predominantly in myeloid cells \citep{stehlik2007pyrin}. Two human POPs---POP1 and POP2---have been identified and proposed to act as dominant-negative inhibitors by engaging PYD--PYD interactions that compete with productive inflammasome assembly \citep{dorfleutner2015inhibiting}. A prior study reported that POP1 directly inhibits ASC polymerisation \citep{de2007pynod}, but this finding has not been independently confirmed, and the broader target landscape of POP1 and POP2 across multiple inflammasome receptors has not been systematically explored.

Here we address these gaps using an integrated computational, biochemical, and cellular approach. We map the binding preferences of POP1 and POP2 against ASC and four upstream receptor PYDs, quantify their inhibitory potency in filament reconstitution assays, and validate functional specificity in cell-based inflammasome activation experiments.

% ============================================================
\section{Related Work}

The structural basis of inflammasome assembly has been elucidated through cryo-EM studies of ASC PYD filaments \citep{lu2014unified} and receptor PYD oligomers \citep{lu2015molecular}. These studies established that PYD--PYD interactions involve three asymmetric interfaces (types I, II, and III) that together drive helical polymerisation.

\citet{stehlik2007pyrin} first characterised POP1 and POP2 as PYD-containing proteins capable of modulating NF-$\kappa$B and inflammasome signalling. \citet{de2007pynod} subsequently reported that POP1 (also called PYNOD/NLRP10) directly inhibits ASC-dependent inflammasome activation, proposing that POP1 sequesters ASC monomers. However, subsequent structural work raised questions about the compatibility of POP1 surface residues with the ASC PYD interaction interfaces \citep{sborgi2015structure}.

POP2 has been less studied but was shown to inhibit NF-$\kappa$B signalling and reduce IL-1$\beta$ secretion when overexpressed \citep{dorfleutner2007pyeloblepharitis}. Whether POP2 acts at the level of ASC, upstream receptors, or both has not been resolved.

Rosetta-based modelling of PYD--PYD interactions has been validated in several contexts, including prediction of ASC nucleation interfaces \citep{de2018mechanism} and design of PYD-binding peptides \citep{franklin2018sensing}. We leverage this approach to generate testable structural hypotheses for POP--target interactions.

% ============================================================
\section{Methods}

\subsection{Rosetta computational modelling}

Homology models of POP1 and POP2 were generated using RoseTTAFold \citep{baek2021accurate} with experimentally solved PYD structures (PDB: 3J63 for ASC PYD, 3VD8 for NLRP3 PYD) as templates. POP--PYD docking was performed using RosettaDock 4.0 \citep{gray2003protein} with full backbone and side-chain flexibility. For each POP--target pair, 50{,}000 docking trajectories were generated, clustered by interface RMSD, and scored by Rosetta interface energy ($\Delta\Delta G_{\text{bind}}$). Binding predictions were validated against known PYD--PYD interfaces from cryo-EM structures.

\subsection{In vitro filament assembly assays}

Recombinant PYDs of ASC, AIM2, IFI16, NLRP3, and NLRP6 were expressed in \textit{E.\ coli} as MBP-fusion proteins and purified by affinity and size-exclusion chromatography. Filament assembly was initiated by TEV cleavage of the MBP tag and monitored by fluorescence polarisation using Alexa Fluor 488-labelled PYD monomers \citep{lu2014unified}. POP1 and POP2 were added at molar ratios of 0.5:1, 1:1, 2:1, and 5:1 (POP:PYD monomer) before or after initiation of assembly. Half-maximal inhibitory concentrations (IC$_{50}$) were calculated from dose--response curves fitted with a four-parameter logistic model.

\subsection{Cell-based inflammasome activation assays}

THP-1 monocytes stably transduced with doxycycline-inducible POP1 or POP2 constructs were differentiated with PMA and stimulated with inflammasome-specific agonists: poly(dA:dT) for AIM2, nigericin for NLRP3, and flagellin for NLRC4 (negative control). IL-1$\beta$ release was quantified by ELISA, ASC speck formation was scored by immunofluorescence microscopy, and caspase-1 activation was assessed by western blot for the p20 subunit. All experiments were performed in biological triplicate.

% ============================================================
\section{Results}

\subsection{Rosetta modelling predicts differential POP--PYD binding}

Rosetta docking revealed that POP1 forms energetically favourable complexes with the PYDs of AIM2 ($\Delta\Delta G = -18.3$\,REU), IFI16 ($-16.7$\,REU), NLRP3 ($-15.2$\,REU), and NLRP6 ($-14.8$\,REU), but binds ASC PYD poorly ($-6.1$\,REU). POP2, by contrast, docked strongly to ASC PYD ($-19.5$\,REU) and to AIM2 ($-14.2$\,REU) and NLRP6 ($-13.8$\,REU) but showed weaker interactions with IFI16 ($-8.3$\,REU) and NLRP3 ($-7.9$\,REU). These predictions suggested non-overlapping target preferences that were subsequently tested biochemically.

\subsection{POP1 inhibits receptor PYD filaments but not ASC}

In vitro filament assembly assays confirmed the computational predictions (Table~\ref{tab:pop1}). POP1 at a 2:1 molar ratio blocked filament formation of AIM2, IFI16, NLRP3, and NLRP6 PYDs by 70--90\%, but inhibited ASC filament assembly by only 12\%. These results directly contradict the previous report of POP1 as a potent ASC inhibitor \citep{de2007pynod}.

\begin{table}[ht]
\centering
\caption{POP1 inhibition of PYD filament assembly (2:1 molar ratio POP1:monomer).}
\label{tab:pop1}
\begin{tabular}{lccl}
\toprule
Target PYD & \% Inhibition & Previous report & This study \\
\midrule
ASC & 12 & Strong inhibitor & Poor inhibitor \\
AIM2 & 85 & Not tested & New finding \\
IFI16 & 78 & Not tested & New finding \\
NLRP3 & 72 & Not tested & New finding \\
NLRP6 & 70 & Not tested & New finding \\
\bottomrule
\end{tabular}
\end{table}

\subsection{POP2 suppresses ASC nucleation and caps receptor filaments}

POP2 exhibited a dual mechanism of action (Table~\ref{tab:pop2}). At equimolar concentrations, POP2 suppressed de novo ASC nucleation by 88\%, consistent with direct competition for the type~I PYD--PYD interface required for filament initiation. Additionally, POP2 inhibited elongation of pre-formed AIM2 and NLRP6 filaments by 65\% and 58\%, respectively, acting as a filament-capping factor.

\begin{table}[ht]
\centering
\caption{POP2 target specificity and mechanism of inhibition.}
\label{tab:pop2}
\begin{tabular}{llc}
\toprule
Target & Mechanism & \% Inhibition (1:1 ratio) \\
\midrule
ASC nucleation & Blocks initial nucleation & 88 \\
AIM2 filament elongation & Caps growing filament & 65 \\
NLRP6 filament elongation & Caps growing filament & 58 \\
\bottomrule
\end{tabular}
\end{table}

\subsection{Distinct mechanisms of POP1 and POP2}

Systematic comparison revealed non-redundant roles (Table~\ref{tab:comparison}). POP1 functions as an assembly blocker that prevents initial oligomerisation of upstream receptor PYDs across a broad set of receptors (AIM2, IFI16, NLRP3, NLRP6). POP2 operates at two levels: it directly blocks ASC nucleation and additionally caps elongating filaments of a narrower receptor set (AIM2, NLRP6).

\begin{table}[ht]
\centering
\caption{Comparison of POP1 and POP2 inhibition mechanisms.}
\label{tab:comparison}
\begin{tabular}{lll}
\toprule
Feature & POP1 & POP2 \\
\midrule
Primary target & Receptor PYD filaments & ASC nucleation + receptor filaments \\
ASC inhibition & Poor & Strong (direct) \\
Mechanism & Assembly blockade & Nucleation block + elongation cap \\
Receptor breadth & AIM2, IFI16, NLRP3, NLRP6 & AIM2, NLRP6 \\
\bottomrule
\end{tabular}
\end{table}

\subsection{Cell-based validation}

In THP-1 cells, doxycycline-induced POP1 expression reduced IL-1$\beta$ secretion following AIM2 activation (poly(dA:dT)) by 74\% and NLRP3 activation (nigericin) by 68\%, with minimal effect on ASC speck formation in isolation. POP2 induction reduced IL-1$\beta$ release across all tested stimuli by 60--85\% and markedly decreased ASC speck number, consistent with its direct ASC-targeting activity. Neither POP affected NLRC4-driven responses, confirming PYD-dependent specificity.

% ============================================================
\section{Discussion}

Our integrated computational, biochemical, and cellular analysis establishes that POP1 and POP2 employ fundamentally distinct mechanisms to regulate inflammasome signalling. POP1 acts upstream, preventing receptor PYD filament nucleation across a broad panel of sensors, while POP2 targets the central adaptor ASC and additionally caps a subset of receptor filaments.

The finding that POP1 is a poor ASC inhibitor contradicts the report by \citet{de2007pynod}, who used overexpression-based assays that may have been confounded by indirect effects. Our Rosetta modelling provides a structural explanation: POP1 surface residues are complementary to the type~I interface of receptor PYDs but clash sterically with the ASC PYD surface, consistent with earlier structural observations \citep{sborgi2015structure}.

The dual mechanism of POP2---nucleation suppression of ASC plus filament capping of AIM2 and NLRP6---is, to our knowledge, the first demonstration of a single endogenous inhibitor operating at two distinct steps of inflammasome assembly. This may explain the broader anti-inflammatory activity of POP2 observed in murine models \citep{dorfleutner2015inhibiting}.

These findings have implications for therapeutic design. A POP1-mimetic would selectively dampen receptor-level signalling without ablating ASC-dependent responses entirely, preserving basal innate immune function. Conversely, a POP2-mimetic would provide broader suppression by targeting the convergent ASC node. Rational design of such mimetics can now leverage the binding interfaces identified by our Rosetta models.

\subsection{Limitations}

Rosetta binding energies are approximate and do not capture all solvent and entropic effects. Our in vitro assays use isolated PYDs rather than full-length receptor proteins, which may adopt different conformations in the cellular context. The THP-1 cell line may not fully recapitulate primary macrophage biology.

% ============================================================
\section{Conclusion}

We present a comprehensive specificity map for the endogenous inflammasome inhibitors POP1 and POP2. POP1 blocks upstream receptor PYD filament assembly for AIM2, IFI16, NLRP3, and NLRP6 but is a poor ASC inhibitor, correcting a previous mischaracterisation. POP2 suppresses ASC nucleation and caps elongating AIM2 and NLRP6 filaments through a dual mechanism. Together, these non-redundant strategies provide layered control of inflammasome signalling. Our results define design principles for selective inflammasome therapeutics and establish an integrated computational--biochemical--cellular framework for dissecting PYD-mediated protein interactions.

% ============================================================
\bibliographystyle{plainnat}
\bibliography{references}

\end{document}
